\section{深度学习不是万能锤子}

\subsection{是否真的需要深度学习？}

讲者提出了一个在深度学习课堂上看似``大逆不道''的问题：\textbf{Do you really want to use deep learning?}

他的核心观点并非否定深度学习，而是提醒学生：

\begin{knowledgebox}{工具选择的智慧}
深度学习是工具箱中非常强大的一个工具，但它不应该是你每次都最先拿出来的那个。计算机科学和传统 AI（如今已融入一般计算机科学）提供了大量替代方案。使用更传统的方法可能意味着：
\begin{itemize}
    \item 更低的风险——行为更可预测、可解释
    \item 更少的数据需求——不需要海量训练数据
    \item 更低的成本——不需要昂贵的 GPU 集群
    \item 更好的可控性——基于规则的系统不会``幻觉''
\end{itemize}
但深度学习在感知任务（图像、语音、自然语言）上的表现往往远超传统方法，因此关键在于\textbf{根据项目需求选择合适的工具}。
\end{knowledgebox}

\subsection{风险无法归零}

讲者强调：\textbf{所有工程都有风险}。无论使用深度学习还是传统方法，风险永远不可能为零。但可以通过以下策略降低风险：

\begin{itemize}
    \item 在能使用传统方法的场景下避免使用深度学习
    \item 在必须使用深度学习的场景下加强测试和评估
    \item 在系统输出端添加额外的安全过滤层
    \item 持续监控系统在生产环境中的表现
\end{itemize}

\subsection{本章小结}

深度学习是强大的工具，但不是唯一的工具。负责任的 AI 开发者应该先评估项目是否真的需要深度学习，然后在风险和能力之间做出权衡。传统方法在可预测性和可控性方面往往具有优势。

